\section{Tests and Evaluation}
\label{sec:evaluation}

\subsection{Testing Strategy}
\label{sec:strategy}

The transformation is tested at four levels. Each level answers a different
question, and the build runs all four on every change.

\begin{enumerate}
  \item \textbf{Parity on the dependency edges.} For each example, the
    pipeline resolves the project and exports the resolved dependency edges as
    canonical JSON. A parity test compares the file byte for byte with the
    committed oracle \texttt{expected/dependencies.json}, which the production
    implementation is tested against as well. The edges are the part of a
    resolved model that both implementations must agree on before either
    generates files. The parity test reads only the files the pipeline wrote,
    so it tests the pipeline as a user runs it.
  \item \textbf{Equality with a hand-written target model.} Task~1 included a
    target model of \texttt{minimal}, written by hand from the specification
    before the transformation existed. A test compares the transformation's
    output with it using EMF's structural equality, \texttt{EcoreUtil.equals}.
    This covers every value of the model, the generated-file entries included.
  \item \textbf{A cross-check against the production implementation.} The
    production implementation writes a resolved document per application,
    \texttt{expected/resolved.json}. Its shape differs from the target
    metamodel, so it is not an oracle for this implementation. During
    development, a script compared every value that both forms record
    (Section~\ref{sec:crosscheck}).
  \item \textbf{Tests of the transformation's own interface.} JUnit tests run
    the transformation on models built in code: the smallest project that
    resolves, a project the input does not contain, and an application whose
    processes need two kinds of workload, which must stop the run. Other tests
    check that each grammar of the pinned inputs reads its document, and that
    a rejected input set is never resolved.
\end{enumerate}

The Java code around the transformation (the pipeline, the runner, the
black-box library and the exporter) is held to 100\% line and branch coverage,
and to a mutation score of 100\%: every change a mutation tool makes to it must
fail a test. The QVTo code is not measured by these tools; the first three
levels test it.

\subsection{Results}
\label{sec:results}

Table~\ref{tab:results} lists the seven examples. All seven resolve, and for
each the exported edges equal the oracle. The output of \texttt{minimal} equals
the hand-written model.

\begin{table}[!htbp]
\centering
\small
\begin{tabular}{lrrrrrrl}
\hline
\textbf{Example} & \textbf{Apps} & \textbf{Proc.} & \textbf{Edges} &
  \textbf{Secrets} & \textbf{Volumes} & \textbf{Env.} & \textbf{Edges equal oracle} \\
\hline
\texttt{minimal}       & 1 & 1 & 0 & 0 & 0 & 6  & yes \\
\texttt{auth}          & 1 & 2 & 4 & 3 & 0 & 21 & yes \\
\texttt{data}          & 3 & 3 & 0 & 1 & 3 & 3  & yes \\
\texttt{delivery}      & 2 & 2 & 0 & 0 & 0 & 4  & yes \\
\texttt{edge}          & 2 & 2 & 0 & 0 & 0 & 0  & yes \\
\texttt{observability} & 2 & 2 & 0 & 0 & 0 & 0  & yes \\
\texttt{secrets}       & 1 & 1 & 0 & 0 & 0 & 0  & yes \\
\hline
\end{tabular}
\caption{The Resolved Examples. Columns Count Applications, Processes,
Dependency Edges, Secrets, Volumes and Environment Variables in Each Resolved
Model.}
\label{tab:results}
\end{table}

The examples are small, but together they use every rule of
Table~\ref{tab:rules}. \texttt{auth} has two processes switched together behind
one release gate, four edges (three to processes in \texttt{data} and one to
the external mail server), two key-value secrets and a transit key, a database
migration with a proof, and dependency, exposure and identity placeholders. \texttt{data}
has three stateful processes, one volume of each durability class, a sidecar,
a configuration file, and a secret placeholder that becomes a reference to a
key of a secret.
Producing the output of all examples, the seven resolutions included, takes
about seven seconds in the build.

Four of the seven oracles, those of \texttt{delivery}, \texttt{edge},
\texttt{observability} and \texttt{secrets}, did not exist before this task.
They were generated by the production implementation, reviewed, committed, and
added to its test suite, so both implementations are tested against them.

\subsubsection{Cross-Check Against the Production Implementation}
\label{sec:crosscheck}

For \texttt{auth}, \texttt{data}, \texttt{delivery} and \texttt{minimal}, the
script compared these values of every application and process: alert class,
secret store, migration, release gate deadline and members, identity, image,
user, group, upgrade mode, switchover, deadline, replica count, memory, CPU,
hardening, token mounting, eligible machines, bound machine, environment
variables with their secret references, writable paths, secrets with their
paths and restart targets, volumes with their backup plans, configuration file
names, sidecars, and the startup probe. All values were equal.

Two values are computed by both implementations but differ by design. The
revision and the input digests are hashes over each implementation's own form
of the model, and the two forms differ. Each implementation is consistent with
itself, and only the production implementation's digests reach the cluster.
The configuration file's name also contains a digest, and there the
transformation hashes the text as the production implementation does, because
the name appears in the generated files.

\subsubsection{Defects Found by Testing}
\label{sec:defects}

The tests found defects that the examples alone would not have shown. Writing
the interface tests showed that a process that runs once, a \emph{job}, was
counted as a blue/green process, so a project of only jobs stopped with a
release-gate error. A review against the production implementation found that
two routes to one port admitted the same proxy twice, that two exposures on one
entry point produced two index files with the same path, that the startup probe
had no timeout, and that a memory amount written in bytes was misread. Each was
fixed, and each has a test or a stop that would catch it again.

\subsection{Limitations}
\label{sec:limitations}

\begin{itemize}
  \item \textbf{The \texttt{knowledge} example is not resolved.} The production
    implementation rejects it, because its images are not in the lock, and its
    committed edge oracle names processes that the \texttt{data} example has
    since renamed. It will be added when the production implementation
    resolves it.
  \item \textbf{Some generated files have no entry yet.} The target metamodel
    describes namespaces, indexes, workloads, switch definitions, service
    accounts, metrics definitions, network policies and routes. It does not
    yet describe volume claims, backup jobs, configuration maps, services,
    disruption budgets, migration jobs or the Vault and secret-operator
    objects. Task~3 adds them with the templates, so the generated-file
    entries of \texttt{auth} and \texttt{data} are incomplete today. Their
    resolved values, which the files will be generated from, are complete.
  \item \textbf{Two places are less detailed than the production model.} The
    order between reconcile units is computed per project, as the union over
    its applications, because Flux orders whole units. A network peer does not
    record which rule admitted it, so two rules that admit the same peer become
    one entry, which admits the same traffic.
  \item \textbf{The workload strategy of the delivery tools differs.} For a
    process that switches by rolling update, the transformation derives a
    rolling update, while the production implementation's file generator
    currently writes a recreate strategy. The specification describes a
    rolling update. The difference will be resolved when the delivery tools
    are generated.
  \item \textbf{Checks across documents are stops, not diagnostics.} An image
    missing from the lock stops the run with one message. A checker would list
    every such problem with its location in the document, as the constraints
    of Task~1 do.
  \item \textbf{The Eclipse launch is a convenience.} The build runs the
    transformation through the standalone executor. The launch configuration
    and the library registration for Eclipse are not run by the build.
\end{itemize}

\subsection{Reflection}
\label{sec:reflection}

QVTo suited this transformation. Most rules are queries over the whole
cluster, which OCL expresses well, and the explicit trace made every reference
between created elements visible in the code. Three properties of the language
cost time: \texttt{resolve}, \texttt{from} and \texttt{access} are keywords and
cannot be used as names, a multi-statement \texttt{if} needs the brace form,
and inside an object expression \texttt{result} still refers to the enclosing
mapping's result, not to the object being created. The compiler reports each of these clearly.

Testing against another implementation's oracles worked well. The dependency
edges were a narrow enough contract to agree on, and the cross-check of the
other values found no difference that the edges had hidden. The hand-written
target model of Task~1 was also useful: it fixed what the transformation had to
produce for \texttt{minimal} before any rule was written. The main cost of the
approach is that every change to the target metamodel changes that model too.
